\documentclass[11pt,a4paper]{article}

% ── Codificación y fuentes ────────────────────────────────────────────────────
\usepackage[utf8]{inputenc}
\usepackage[T1]{fontenc}
\usepackage{lmodern}
\usepackage[spanish]{babel}

% ── Geometría y espaciado ─────────────────────────────────────────────────────
\usepackage[top=2.5cm, bottom=2.5cm, left=2.8cm, right=2.8cm]{geometry}
\usepackage{setspace}
\setstretch{1.15}
\usepackage{parskip}

% ── Tablas ────────────────────────────────────────────────────────────────────
\usepackage{booktabs}
\usepackage{tabularx}
\usepackage{array}
\usepackage{longtable}
\usepackage{enumitem}

% ── Colores y cajas ───────────────────────────────────────────────────────────
\usepackage[dvipsnames]{xcolor}
\usepackage{tcolorbox}
\tcbuselibrary{skins, breakable}

% ── Cabeceras y pies de página ────────────────────────────────────────────────
\usepackage{fancyhdr}
\usepackage{lastpage}
\pagestyle{fancy}
\fancyhf{}
\fancyhead[L]{\small\textbf{Informe de Evaluación} --- Electrónica}
\fancyhead[R]{\small Sebastian Pulido}
\fancyfoot[C]{\small Página \thepage\ de \pageref{LastPage}}
\renewcommand{\headrulewidth}{0.4pt}

% ── Hiperenlaces ──────────────────────────────────────────────────────────────
\usepackage[colorlinks=true, linkcolor=NavyBlue, urlcolor=NavyBlue]{hyperref}

% ── Paleta de colores personalizados ─────────────────────────────────────────
\definecolor{desempeno_excelente}{HTML}{1A6B2F}
\definecolor{desempeno_bueno}{HTML}{2E5A9C}
\definecolor{desempeno_suficiente}{HTML}{8B6914}
\definecolor{desempeno_deficiente}{HTML}{C0392B}
\definecolor{desempeno_insuficiente}{HTML}{7B241C}
\definecolor{grisFondo}{HTML}{F5F5F5}
\definecolor{azulTitulo}{HTML}{1B2A6B}
\definecolor{verdeFortaleza}{HTML}{1A6B2F}
\definecolor{naranjaMejora}{HTML}{A04000}

% ── Estilos de cajas ─────────────────────────────────────────────────────────
\tcbset{
  cajaTitulo/.style={
    enhanced, breakable,
    colback=azulTitulo!8, colframe=azulTitulo,
    fonttitle=\bfseries\large, coltitle=white,
    attach boxed title to top left={yshift=-2mm, xshift=4mm},
    boxed title style={colback=azulTitulo},
    top=6pt, bottom=6pt, left=8pt, right=8pt,
  },
  cajaFortaleza/.style={
    enhanced, breakable,
    colback=verdeFortaleza!6, colframe=verdeFortaleza!60,
    fonttitle=\bfseries, coltitle=verdeFortaleza,
    top=4pt, bottom=4pt, left=8pt, right=8pt,
  },
  cajaMejora/.style={
    enhanced, breakable,
    colback=naranjaMejora!6, colframe=naranjaMejora!60,
    fonttitle=\bfseries, coltitle=naranjaMejora,
    top=4pt, bottom=4pt, left=8pt, right=8pt,
  },
  cajaRecomendacion/.style={
    enhanced, breakable,
    colback=grisFondo, colframe=gray!50,
    fonttitle=\bfseries, coltitle=black,
    top=4pt, bottom=4pt, left=8pt, right=8pt,
  },
}

% ─────────────────────────────────────────────────────────────────────────────
\begin{document}

% ══ PORTADA ══════════════════════════════════════════════════════════════════
\begin{center}
  {\Large\bfseries\color{azulTitulo} ELECTRÓNICA --- Evaluación de Exámenes}\\[4pt]
  {\large Informe de Evaluación Preliminar}\\[12pt]
  \rule{\linewidth}{1.2pt}\\[6pt]
  {\LARGE\bfseries Sebastian Pulido}\\[4pt]
  \rule{\linewidth}{0.6pt}
\end{center}

% ══ FICHA TÉCNICA ════════════════════════════════════════════════════════════
\vspace{4pt}
\begin{tcolorbox}[cajaTitulo, title=Datos del informe]
\begin{tabular}{@{}llll@{}}
  \textbf{Estudiante:}  & Sebastian Pulido  &
  \textbf{Fecha:}       & 2026-06-26 \\[4pt]
  \textbf{Archivo PDF:} & \texttt{examen\_Sebastian\_Pulido.pdf}  &
  \textbf{Páginas:}     & 2 \\
\end{tabular}
\end{tcolorbox}

% ══ RESULTADO GLOBAL ═════════════════════════════════════════════════════════
\vspace{6pt}
\begin{center}
  \begin{tcolorbox}[
    enhanced,
    colback=white, colframe=desempeno_suficiente,
    width=0.62\linewidth,
    boxrule=2pt,
    halign=center, valign=center,
    top=8pt, bottom=8pt,
  ]
    \centering
    {\large\bfseries Puntaje sugerido}\\[6pt]
    {\Huge\bfseries\color{desempeno_suficiente} 6.0/10}\\[6pt]
    {\large Nivel de desempeño: \textbf{\color{desempeno_suficiente} Suficiente}}
  \end{tcolorbox}
\end{center}

% ══ RESUMEN GENERAL ══════════════════════════════════════════════════════════
\section*{Resumen general}
El estudiante demuestra una comprensión sólida en la aplicación del Teorema de Máxima Transferencia de Potencia y en el planteamiento de ecuaciones de nodos para supernodos con fuentes dependientes. Sin embargo, presenta dificultades en la interpretación de problemas sin diagramas de circuito explícitos, realiza suposiciones no justificadas y comete un error de signo en la aplicación de la Ley de Voltajes de Kirchhoff (KVL) para una fuente dependiente en el análisis de mallas. La claridad y organización de su trabajo son generalmente buenas.

% ══ FORTALEZAS Y ÁREAS DE MEJORA ════════════════════════════════════════════
\vspace{4pt}
\begin{minipage}[t]{0.48\linewidth}
  \begin{tcolorbox}[cajaFortaleza, title={Fortalezas}]
\begin{itemize}[leftmargin=1.5em, itemsep=2pt]
  \item Correcta aplicación del Teorema de Máxima Transferencia de Potencia, incluyendo la fórmula y el cálculo.
  \item Habilidad para identificar y plantear ecuaciones para supernodos y fuentes dependientes en análisis nodal a partir de una descripción textual.
  \item Buena capacidad de manipulación algebraica para simplificar ecuaciones.
  \item Presentación legible y organizada de las respuestas.
\end{itemize}
  \end{tcolorbox}
\end{minipage}
\hfill
\begin{minipage}[t]{0.48\linewidth}
  \begin{tcolorbox}[cajaMejora, title={Areas de mejora}]
\begin{itemize}[leftmargin=1.5em, itemsep=2pt]
  \item Interpretación rigurosa de los enunciados de los problemas, especialmente cuando no se proporcionan diagramas de circuito.
  \item Justificación de las suposiciones realizadas durante el análisis de circuitos.
  \item Aplicación precisa de la Ley de Voltajes de Kirchhoff (KVL), prestando especial atención a la polaridad de las fuentes de tensión, particularmente las dependientes.
  \item Completitud en la presentación de todas las ecuaciones necesarias para resolver un sistema (ej. todas las ecuaciones de malla).
\end{itemize}
  \end{tcolorbox}
\end{minipage}

% ══ OBSERVACIONES POR PREGUNTA ═══════════════════════════════════════════════
\section*{Observaciones por pregunta}

\begin{longtable}{>{\bfseries}p{2.8cm} p{1.6cm} p{7.0cm} p{2.8cm}}
  \toprule
  \textbf{Pregunta} & \textbf{Puntaje} & \textbf{Evaluación} & \textbf{Errores detectados} \\
  \midrule
  \endhead
  \bottomrule
  \endfoot
    \textbf{Pregunta 1} & 0.5/2.0 & La respuesta es confusa y contradictoria. El estudiante presenta V\_x=5V y V\_ab=-2V como datos, y luego calcula V\_ab=3V. Sin el diagrama del circuito original, es imposible verificar la validez de estas variables o la operación realizada. La contradicción de V\_ab siendo -2V y luego 3V es un error fundamental de interpretación o aplicación. & \textit{Contradicción en el valor de V\_ab (dado como -2V y calculado como 3V).; Falta de contexto (diagrama de circuito) que hace imposible verificar la interpretación de las variables V\_x y V\_ab.} \\
    \midrule
    \textbf{Pregunta 2} & 0.5/2.0 & El estudiante asume 'No hay corriente: I\_T = 0A' y 'No hay caída de tensión: V\_ab = 0V' sin justificación, lo cual es contradictorio (si no hay corriente, no hay caída de tensión en resistencias, pero V\_ab=0V no es necesariamente el V\_th). La determinación de V\_th = 10V podría ser correcta bajo la suposición de circuito abierto, pero la R\_th = 2$\Omega$ es muy cuestionable dado los valores de resistencia (3$\Omega$, 2$\Omega$, 3$\Omega$) y sin un diagrama de circuito. Parece una suposición arbitraria o un error de cálculo. & \textit{Suposiciones no justificadas ('No hay corriente', 'No hay caída de tensión').; Contradicción entre las suposiciones (I\_T=0A y V\_ab=0V).; Cálculo de R\_th = 2$\Omega$ parece incorrecto o no justificado a partir de los valores de resistencia dados.} \\
    \midrule
    \textbf{Pregunta 3} & 2.0/2.0 & Excelente respuesta. El estudiante comprende y aplica correctamente el Teorema de Máxima Transferencia de Potencia. La condición para la máxima transferencia de potencia es correctamente enunciada, la fórmula utilizada es la adecuada y los cálculos son precisos, incluyendo la conversión de unidades. & \textit{---} \\
    \midrule
    \textbf{Pregunta 4} & 2.0/2.0 & Muy buena respuesta. El estudiante identifica correctamente el concepto de supernodo y la definición de la fuente dependiente. El planteamiento de la ecuación del supernodo (V1 - V2 = 4Ix) y la ley de Ohm para Ix (Ix = V2 / 2) son correctos, y la sustitución y simplificación algebraica (V1 = 3V2) se realizan sin errores. La interpretación de la descripción textual del problema es acertada. & \textit{---} \\
    \midrule
    \textbf{Pregunta 5} & 1.0/2.0 & El estudiante identifica correctamente la resistencia común entre mallas y el principio de KVL. Sin embargo, al aplicar KVL en la Malla 2, comete un error de signo al incluir la fuente de tensión dependiente (3I1). Según el diagrama que él mismo dibuja (polaridad + arriba, - abajo), al recorrer la malla en sentido horario, se pasa de - a +, por lo que el término debería ser -3I1, no +3I1. Además, solo se presenta la ecuación para la Malla 2, faltando la ecuación para la Malla 1 para un análisis completo. & \textit{Error de signo en la aplicación de KVL para la fuente de tensión dependiente (3I1) en la Malla 2.; Incompletitud: Solo se presenta una de las dos ecuaciones de malla necesarias para resolver el circuito.} \\
    \midrule
\end{longtable}

% ══ ERRORES DETECTADOS ═══════════════════════════════════════════════════════
\section*{Análisis de errores}

\subsection*{Errores conceptuales}
\begin{itemize}[leftmargin=1.5em, itemsep=2pt]
  \item Confusión o uso contradictorio de variables en la Pregunta 1, lo que sugiere una falta de comprensión clara del escenario planteado.
  \item Suposiciones no justificadas en la Pregunta 2 que llevan a una solución potencialmente incorrecta para R\_th.
  \item Error en la aplicación de la polaridad de una fuente de tensión dependiente al aplicar KVL en la Pregunta 5.
\end{itemize}

\subsection*{Errores procedimentales}
\begin{itemize}[leftmargin=1.5em, itemsep=2pt]
  \item Cálculo de R\_th en la Pregunta 2 que no se deriva lógicamente de los valores de resistencia dados sin un diagrama.
  \item Incompletitud en la Pregunta 5 al no proporcionar todas las ecuaciones de malla requeridas para el análisis.
\end{itemize}

% ══ RECOMENDACIONES AL ESTUDIANTE ════════════════════════════════════════════
\section*{Retroalimentación para el estudiante}
\begin{tcolorbox}[cajaRecomendacion, title={Dirigido a: Sebastian Pulido}]
Para mejorar tu desempeño, te sugiero que siempre intentes dibujar el circuito cuando no se te proporciona, basándote en la descripción textual. Esto te ayudará a visualizar las conexiones y polaridades. Presta especial atención a la aplicación de las leyes de Kirchhoff, revisando cuidadosamente los signos de las fuentes de tensión y las caídas de tensión en las resistencias, especialmente con fuentes dependientes. Además, asegúrate de justificar cualquier suposición que hagas y de presentar un conjunto completo de ecuaciones cuando sea necesario para resolver un problema.
\end{tcolorbox}

% ══ NOTA PARA EL PROFESOR ════════════════════════════════════════════════════
\section*{Nota para el profesor}
\begin{tcolorbox}[
  enhanced, breakable,
  colback=yellow!8, colframe=yellow!60!black,
  fonttitle=\bfseries, coltitle=black,
  title={Observaciones para revisión manual},
  top=4pt, bottom=4pt, left=8pt, right=8pt,
]
Las Preguntas 1 y 2 son difíciles de evaluar con total certeza debido a la ausencia de los diagramas de circuito originales. La interpretación del estudiante podría estar basada en un circuito implícito que no se conoce. Sería útil revisar si las variables y condiciones dadas en estas preguntas corresponden a un circuito específico que los estudiantes debían conocer o inferir. La interpretación textual para las Preguntas 4 y 5 parece razonable por parte del estudiante, a pesar del error de signo en KVL en la Pregunta 5.
\end{tcolorbox}

% ══ PIE DE INFORME ════════════════════════════════════════════════════════════
\vspace{12pt}
\begin{center}
  \footnotesize\color{gray}
  Informe generado automáticamente por el sistema de evaluación con IA (gemini-2.5-flash) y soporte LaTex. \\
  Este documento es una evaluación \textbf{preliminar} para ser revisado por el profesor antes de ser entregado al 
  estudiante. \\
  Generado el 2026-06-26 \textbullet\ ELECTRÓNICA --- Sistema de Evaluación Académica.
  \end{center}

\end{document}
